% 2.4 =============== 最大值原理与值函数性质 ===============

\section{值函数的存在性与性质} \label{sec-value-func-property}

值函数连续性和相关性质的证明相对较为复杂，我们通过引入最优性原理（Principal of Optimality）进而展开。


% 2.4.1 ===== 最优性原理 =====
\subsection{最优性原理：序列解与递归解的关系}

首先，将序列问题（Sequential Problem，$\mathbf{SP}$）写为一般化形式：
\begin{equation*}
    \begin{aligned}
        \sup _{\left\{x_{t+1}\right\}_{t=0}^{\infty}} & \sum_{t=0}^{\infty} \beta^t F\left(x_t, x_{t+1}\right) \\
        \text { s.t. } & x_{t+1} \in \Gamma\left(x_t\right), \quad \forall t \\
        & x_0 \in X \quad \text { 给定 }
    \end{aligned}
    \tag{SP} \label{problem-sp}
\end{equation*}
其中，$X$是$x$所有可能取值集合；$\Gamma : X \rightarrow X$是描述可行性的对应（Correspondence），即
给定当前状态$x_t$，下期$x_{t+1}$的所有可行取值；$F$是当期回报函数。我们称从$x_0$开始，
任意可行的序列$\{x_{t+1}\}_{t=0}^{\infty}$为一个可行计划（Feasible Plans）。记所有可行计划构成的集合为：
\begin{equation*}
    \pi\left(x_0\right)=\left\{\left\{x_{t+1}\right\}_{t=0}^{\infty}: ~ x_{t+1} \in \Gamma\left(x_t\right), 
        ~ t=0,1,2, \ldots\right\}
\end{equation*}
其元素是一个序列，记为$\boldsymbol{x}\in\pi\left(x_0\right)$。

其次，我们写出与序列问题$SP$相对应的递归问题$R$，或称为泛函方程（Functional Equation，$FE$）：
\footnote{在前面我们研究一般化的贝尔曼方程时使用的记号是$\mathcal{F}(x,x^{\prime})$，因为我们使用$h(x,y)$将控制变量进行了替换，
这里并没有本质区别，只是记号的差异。}
\begin{equation}
    \begin{aligned}
        v(x)=\sup _{y \in \Gamma(x)}[F(x, y)+\beta v(y)] \quad \forall x \in X
    \end{aligned}
    \tag{FE} \label{problem-fe}
\end{equation}
其中$x$是状态变量，$y$是控制变量。

那么，$SP$与$FE$是否等价？更一般的想法是：
\begin{itemize}
    \item 问题$FE$的解（函数）$v$在$x=x_0$处的取值达到$SP$的上界；
    \item 序列$\{x_{t+1}\}_{t=0}^{\infty}$达到问题$SP$上界，当且仅当
        $v\left(x_t\right)=F\left(x_t, x_{t+1}\right)+\beta v\left(x_{t+1}\right), \quad t=0,1,2, \ldots$
\end{itemize}
贝尔曼将这一想法称为\textbf{最优性原理}。当然，这一结果的成立需要一些假设。
\begin{assumption} \label{as-nonempty-plan}
    集合$\Gamma\left(x\right)$对任意$x\in X$非空
\end{assumption}
\begin{assumption} \label{as-lim-exist}
    对任意$\left\{x_0, x_1, \ldots\right\} \in \Pi\left(x_0\right), x_0 \in X$，
    目标极限$\lim\limits_{n\rightarrow \infty}\sum_{t=0}^{n}\beta^t F\left(x_t,x_{t+1}\right)$存在（可能为无穷）
\end{assumption}
不难看出，\cref{as-lim-exist}成立的一个充分条件是：$F\left(x,y\right)$有界（bounded）且$0<\beta<1$。
对任意$n=0,1,\dots$，定义$u_n: \pi\left(x_0\right)\rightarrow \mathbb{R}$：
$$
u_n(\boldsymbol{x}) \equiv \sum_{t=0}^n \beta^t F\left(x_t, x_{t+1}\right), 
    \quad \boldsymbol{x}=\{x_0,x_1,\dots\} \in \pi\left(x_0\right)
$$
由\cref{as-lim-exist}，可定义极限：
\begin{equation*}
    u\left(\boldsymbol{x}\right) \equiv \lim _{n\rightarrow\infty}u_n\left(\boldsymbol{x}\right)
\end{equation*}
如果\cref{as-nonempty-plan}和\cref{as-lim-exist}均成立，那么序列问题$SP$是良定义的。因此我们可以定义\textbf{上确界函数}：
\begin{equation*}
    v^*\left(x_0\right)\equiv \sup _{\boldsymbol{x}\in\pi\left(x_0\right)} u\left(\boldsymbol{x}\right)
\end{equation*}
这样，$v^*\left(x_0\right)$是$SP$问题的上确界。根据上确界定义，其满足：
\begin{equation*}
    \begin{aligned}
        &\forall \epsilon > 0, ~\forall \boldsymbol{x}\in \pi\left(x_0\right): 
            ~ v^*\left(x_0\right)\geq u\left(\boldsymbol{x}\right) \\
        &\forall \epsilon > 0, ~\exists \boldsymbol{x}\in\pi\left(x_0\right): 
            ~ v^*\left(x_0\right)\leq u\left(\boldsymbol{x}\right)+\epsilon
    \end{aligned}
\end{equation*}

我们希望研究$SP$问题的\textbf{上确界函数}$v^*$和$FE$问题的解$v$的关系。
需要注意：$v^*$是唯一定义的，但$FE$问题的解$v$可能有很多个。简单起见，以下不考虑无穷大的情形，
即考虑$\left|v^*\left(x_0\right)\right| < \infty$。那么，$v^*$是$FE$问题的解等价于如下条件成立：
\begin{equation*}
    \begin{aligned}
        &\forall y \in \Gamma\left(x_0\right):
            ~ v^*\left(x_0\right)\geq F\left(x_0,y\right) + \beta v^*\left(y\right) \\
        &\forall \epsilon > 0, ~\exists y\in \Gamma\left(x_0\right):
            ~ v^*\left(x_0\right)\leq F\left(x_0,y\right) + \beta v^*\left(y\right) + \epsilon, 
    \end{aligned} 
\end{equation*}

为了证明$v^*$满足$FE$，引入如下引理：
\begin{lemma}
    令$X, \Gamma, F, \beta$符合\cref{as-lim-exist}，
    则对$\forall x_0 \in X, \forall \left(x_0,x_1,\dots\right)=\boldsymbol{x}\in\pi\left(x_0\right)$，成立：
    \begin{equation*}
        u\left(\boldsymbol{x}\right)=
            F\left(x_0,x_1\right)+\beta u\left(\boldsymbol{x}^{\prime}\right), \quad \boldsymbol{x}^{\prime}
                =\left(x_1,x_2,\dots\right)
    \end{equation*}
\end{lemma}
\begin{theorem}[SP-FE] \label{th-solution-sp-fe}
    令$X, \Gamma, F, \beta$符合\cref{as-nonempty-plan}和\cref{as-lim-exist}，则上确界函数$v^*$满足$FE$。
\end{theorem}
\begin{proof}
    前面提到，上确界函数$v^*$是$FE$问题的解等价于成立如下两个条件：
    \begin{equation*}
        \begin{aligned}
            &\forall y \in \Gamma\left(x_0\right): ~v^*\left(x_0\right)\geq F\left(x_0,y\right) 
                + \beta v^*\left(y\right)\\
            &\forall \epsilon >0, ~\exists y\in \Gamma\left(x_0\right): 
                ~v^*\left(x_0\right)\leq F\left(x_0,y\right) + \beta v^*\left(y\right) + \epsilon
        \end{aligned}
    \end{equation*}

    先证明第一个条件：
    \begin{equation*}
        \forall y \in \Gamma\left(x_0\right):~ v^*\left(x_0\right)\geq F\left(x_0,y\right) + \beta v^*\left(y\right)
    \end{equation*}
    由$v^*\left(x_0\right)\equiv \sup _{\boldsymbol{x}\in\pi\left(x_0\right)} u\left(\boldsymbol{x}\right)$，
    根据上确界的定义，成立如下两条：
    \begin{equation*}
        \begin{aligned}
            &\forall \epsilon > 0, \forall \boldsymbol{x}\in \pi\left(x_0\right): 
                \quad v^*\left(x_0\right)\geq u\left(\boldsymbol{x}\right) \\
            &\forall \epsilon > 0, \exists \boldsymbol{x}\in\pi\left(x_0\right): 
                \quad v^*\left(x_0\right)\leq u\left(\boldsymbol{x}\right)+\epsilon
        \end{aligned}
    \end{equation*}
    根据后者，任取$x_1 \in \Gamma\left(x_0\right)$，任取$\epsilon > 0$：
    \begin{equation*}
        \exists \boldsymbol{x}^{\prime}=\left(x_1,x_2,\dots\right) \in \pi\left(x_1\right), 
            \quad u\left(\boldsymbol{x}^{\prime}\right)\geq v^*\left(x_1\right)-\epsilon
    \end{equation*}
    又由于$\boldsymbol{x}=\left(x_0,x_1,\dots\right)\in \pi\left(x_0\right)$，根据前者和上述引理：
    \begin{equation*}
        v^*\left(x_0\right)\geq u\left(\boldsymbol{x}\right) = 
            F\left(x_0,x_1\right)+\beta u\left(\boldsymbol{x}^{\prime}\right)
                \geq F\left(x_0,x_1\right)+\beta v^*\left(x_1\right)-\beta \epsilon
    \end{equation*}
    由于$\epsilon$任意，可知：
    \begin{equation*}
        \forall y \in \Gamma\left(x_0\right): 
            ~v^*\left(x_0\right)\geq F\left(x_0,y\right) + \beta v^*\left(y\right)
    \end{equation*}

    再证明第二个条件：
    \begin{equation*}
        \forall \epsilon > 0, ~\exists y\in \Gamma\left(x_0\right): 
            ~v^*\left(x_0\right)\leq F\left(x_0,y\right) + \beta v^*\left(y\right) + \epsilon
    \end{equation*}
    任取$x_0\in X$和$\epsilon>0$，根据上确界函数定义的第二条和引理可知：
    \begin{equation*}
        \exists \boldsymbol{x}=\left(x_0,x_1,\dots\right)\in \pi\left(x_0\right), 
            \quad v^*\left(x_0\right)\leq u\left(\boldsymbol{x}\right)+\epsilon 
                = F\left(x_0,x_1\right)+\beta u\left(\boldsymbol{x}^{\prime}\right)+\epsilon
    \end{equation*}
    其中，$\boldsymbol{x}^{\prime}= \left(x_1,x_2,\dots\right)$。从而，结合上确界定义的第一条可知：
    \begin{equation*}
        v^*\left(x_0\right)\leq F\left(x_0,x_1\right)+\beta v^*\left(x_1\right)+\epsilon 
    \end{equation*}
    由$\epsilon$任意性和$x_1\in\Gamma\left(x_0\right)$（因为$x_1$在计划中）：
    \begin{equation*}
        \forall \epsilon > 0, ~\exists y\in \Gamma\left(x_0\right): 
            ~v^*\left(x_0\right)\leq F\left(x_0,y\right) + \beta v^*\left(y\right) + \epsilon
    \end{equation*}
    从而我们证明了$v^*$是$FE$问题的解。
\end{proof}

至此，我们证明了$SP$的解满足$FE$，但反过来是否成立？一般来说不一定，虽然$v^*$是$FE$的解，但$FE$可能存在很多解。
不过下面的定理表明：若$FE$满足特定有界条件，那么$v=v^*$是其唯一解。

\begin{theorem}[Partial Converse: FE-SP] \label{th-solution-fe-sp}
    令$X, \Gamma, F, \beta$符合\cref{as-nonempty-plan}和\cref{as-lim-exist}。如果值函数$v$是$FE$的一个解且满足
    $\lim\limits_{n\rightarrow \infty}\beta^n v\left(x_n\right)=0, ~\forall x_0\in X, 
        \forall \boldsymbol{x}=\left(x_0,x_1,\dots\right)\in \pi\left(x_0\right)$，
    则$v=v^*$。其中：
    \begin{gather*}
        v^*\left(x_0\right) \equiv \sup _{\boldsymbol{x}\in\pi\left(x_0\right)} u\left(\boldsymbol{x}\right) \\
        u\left(\boldsymbol{x}\right) \equiv \lim _{n\rightarrow\infty}u_n\left(\boldsymbol{x}\right) \\
        u_n(\boldsymbol{x}) \equiv \sum_{t=0}^n \beta^t F\left(x_t, x_{t+1}\right), 
            \quad \boldsymbol{x}=\{x_0,x_1,\dots\} \in \pi\left(x_0\right)
    \end{gather*}
\end{theorem}
\begin{proof}
    方便起见仍考虑$v\left(x_0\right)$有限情形。我们此前提到，$v$是$FE$的解等价于如下条件成立：
    \begin{equation*}
        \begin{aligned}
            &\forall y \in \Gamma\left(x_0\right): 
                ~v\left(x_0\right)\geq F\left(x_0,y\right) + \beta v\left(y\right)\\
            &\forall \epsilon > 0, ~\exists y\in \Gamma\left(x_0\right): 
                ~v\left(x_0\right)\leq F\left(x_0,y\right) + \beta v\left(y\right) + \epsilon
        \end{aligned}
    \end{equation*}
    
    我们的目标是证明$FE$的这一解$v$是$SP$的解，即证明
    $v\left(x_0\right)=\sup _{\boldsymbol{x}\in\pi\left(x_0\right)} u\left(\boldsymbol{x}\right)$，
    也即往证$FE$的解$v$符合$SP$中上确界的定义：
    \begin{equation*}
        \begin{aligned}
            &\forall \epsilon > 0, ~\forall \boldsymbol{x}\in \pi\left(x_0\right): 
                ~ v\left(x_0\right)\geq u\left(\boldsymbol{x}\right) \\
            &\forall \epsilon > 0, ~\exists \boldsymbol{x}\in\pi\left(x_0\right): 
                ~ v\left(x_0\right)\leq u\left(\boldsymbol{x}\right)+\epsilon
        \end{aligned}
    \end{equation*}
    注意到对任意$\boldsymbol{x}\in \pi\left(x_0\right)$，成立：
    \begin{equation*}
        \begin{aligned}
            v\left(x_0\right) &\geq F\left(x_0,x_1\right) + \beta v\left(x_1\right) \\
                &\geq  F\left(x_0,x_1\right) + \beta F\left(x_1,x_2\right) + \beta^2v\left(x_2\right) \\
                &\geq \cdots \\
                &\geq u_n\left(\boldsymbol{x}\right)+\beta^{n+1}v\left(x_{n+1}\right)
        \end{aligned}
    \end{equation*}
    取极限并利用
    $\lim\limits_{n\rightarrow \infty}\beta^n v\left(x_n\right)=0, 
        ~ \forall x_0\in X, ~\forall \boldsymbol{x}=\left(x_0,x_1,\dots\right)\in \pi\left(x_0\right)$
    可知上确界定义的前一半成立：
    \begin{equation*}
        \forall \epsilon > 0, ~\forall \boldsymbol{x}\in \pi\left(x_0\right): 
            ~ v\left(x_0\right)\geq u\left(\boldsymbol{x}\right)
    \end{equation*}
    下面证明上确界定义的后一半，即：
    \begin{equation*}
        \forall \epsilon > 0, ~\exists \boldsymbol{x}\in\pi\left(x_0\right): 
            ~ v\left(x_0\right)\leq u\left(\boldsymbol{x}\right)+\epsilon  
    \end{equation*}
    对此，任取一个$\epsilon>0$，从正实数集中选择$\{\delta_t\}_{t=1}^{\infty}$，
    使得$\sum_{t=1}^{\infty}\beta^{t-1}\delta_t\leq \epsilon /2$。由与$v$是$FE$的解，即：
    \begin{equation*}
        \forall \epsilon > 0, ~\exists y\in \Gamma\left(x_0\right): 
            ~v\left(x_0\right)\leq F\left(x_0,y\right) + \beta v\left(y\right) + \epsilon, 
    \end{equation*}
    从而可以取到$x_1\in\Gamma\left(x_0\right), x_2\in\Gamma\left(x_1\right),\cdots$：
    \begin{equation*}
        v\left(x_t\right)\leq 
            F\left(x_t,x_{t+1}\right)+\beta v\left(x_{t+1}\right)+\delta_{t+1},\quad t=0,1,\cdots
    \end{equation*}
    也就是说，$\exists\boldsymbol{x}=\left(x_0,x_1,\cdots\right)\in \pi\left(x_0\right)$：
    \begin{equation*}
        \begin{aligned}
            v\left(x_0\right) &\leq \sum_{t=0}^{n}\beta^t F\left(x_t,x_{t+1}\right) + 
                \beta^{n+1}v\left(x_{n+1}\right) + \left(\delta_1+\cdots+\beta^n \delta_{n+1}\right)   \\
            &\leq u_n\left(\boldsymbol{x}\right)+\beta^{n+1}v\left(x_{n+1}\right)+\epsilon/2, \quad n=1,2,\cdots
        \end{aligned}
    \end{equation*}
    当$n$足够大，由条件$\lim_{n\rightarrow \infty}\beta^n v\left(x_n\right)=0$可知：
    $v\left(x_0\right)\leq u_n\left(\boldsymbol{x}\right)+\epsilon$。
    由$\epsilon$的可知任意性，上确界后半的定义成立：
    \begin{equation*}
        \forall \epsilon > 0, \exists \boldsymbol{x}\in\pi\left(x_0\right): 
            \quad v^*\left(x_0\right)\leq u\left(\boldsymbol{x}\right)+\epsilon
    \end{equation*}
    从而：
    \begin{equation*}
        v\left(x_0\right)=\sup _{\boldsymbol{x}\in\pi\left(x_0\right)} u\left(\boldsymbol{x}\right)
    \end{equation*}
    也即$FE$的解$v$是$SP$的解。
\end{proof}

至此，我们可以对序列问题（$SP$）和递归（泛函）问题（$FE$）解的关系进行一个总结。方便起见，再次写下这两个一般化的问题：
\begin{itemize}
    \item $SP$：
    \begin{equation*}
        \sup _{\left\{x_{t+1}\right\}_{t=0}^{\infty}} \sum_{t=0}^{\infty} \beta^t F\left(x_t, x_{t+1}\right),
            \text { s.t. } ~x_{t+1} \in\Gamma\left(x_t\right), ~\forall t, ~x_0 \in X \text { 给定 }
    \end{equation*}
    \item $FE$：
    \begin{equation*}
        v(x)=\sup _{y \in \Gamma(x)}[F(x, y)+\beta v(y)], ~\forall x \in X
    \end{equation*}
\end{itemize}

在\cref{as-nonempty-plan}和\cref{as-lim-exist}下，我们通过构造极限函数的上确界找到了$SP$的解：
\begin{equation*}
    v^*\left(x_0\right)=\sup _{\boldsymbol{x}\in\pi\left(x_0\right)}u\left(\boldsymbol{x}\right), ~
        u\left(\boldsymbol{x}\right)=\lim_{n\rightarrow\infty}\sum_{t=0}^{n}\beta^t F\left(x_t,x_{t+1}\right)
\end{equation*}
同时证明了它至少也是$FE$的一个解，即$SP\rightarrow FE$。
另一方面，$FE$可能存在很多解$v$，我们证明了其中符合有界条件
$\lim\limits_{n\rightarrow \infty}\beta^n v\left(x_n\right)=0, \quad \forall x_0\in X, 
    \forall \boldsymbol{x}=\left(x_0,x_1,\dots\right)\in \pi\left(x_0\right)$
的解正是$SP$的解$v^*$，即部分可逆。而$FE$的其余解一定违反了这一有界条件。


% 2.4.2 ===== 最优策略的关系 =====
\subsection{最优策略间的关系}

最优性原理（\cref{th-solution-sp-fe}和\cref{th-solution-fe-sp}）给出了序列问题和递归（泛函）问题目标值函数间的关系。
我们希望进一步研究两个问题最优策略（Plan）间的关系。

\begin{definition}[\textbf{最优计划}]
    如果可行计划$\boldsymbol{x}^* = \{x_{t+1}^*\} \in\pi\left(x_0\right)$能达到$SP$问题的上确界，即
    $u\left(\boldsymbol{x}^*\right) \equiv 
        \sum_{t=0}^{\infty}\beta^t F\left(x_t^*,x_{t+1}^*\right) = v^*\left(x_0\right)$，
    则称$\boldsymbol{x}^*$为从$x_0$开始的最优计划（Optimal Plan）。
\end{definition}

下面两个定理给出了最优计划和$FE$问题中满足$v=v^*$的决策规则间的关系。

\begin{theorem} \label{th-policy-sp-fe}
    令$X, \Gamma, F, \beta$符合\cref{as-nonempty-plan}和\cref{as-lim-exist}。
    令$\boldsymbol{x}^*\in \pi\left(x_0\right)$为一个达到$SP$问题上确界的可行计划，则:
    \begin{equation} \label{eq-opt-plan-sp}
        v^*\left(x_t^*\right) = F\left(x_t^*,x_{t+1}^*\right) + \beta v^*\left(x_{t+1}^*\right), ~t=0,1,2,\cdots
    \end{equation}
\end{theorem}

\begin{proof}
    由于$\boldsymbol{x}^*$达到了上确界，因此由\cref{th-solution-sp-fe}，$v^*\left(x_t^*\right)$满足$FE$，
    即对$\forall \boldsymbol{x}\in \pi\left(x_0\right)$：
    \begin{equation*}
        \begin{aligned}
            v^*\left(x_0^*\right) &= u\left(\boldsymbol{x}^*\right) = F\left(x_0,x_1^*\right) 
                + \beta u\left(\boldsymbol{x^{*\prime}}\right) \\
            &\geq u\left(\boldsymbol{x}\right) = F\left(x_0,x_1\right)+\beta u\left(\boldsymbol{x}^{\prime}\right), 
        \end{aligned}
    \end{equation*}
    又由于$\left(x_1^*,x_2,x_3,\cdots\right)\in \pi\left(x_1^*\right)$，
    从而$\left(x_0,x_1^*,x_2,x_3,\cdots\right)\in \pi\left(x_0\right)$,
    进而有：
    \begin{equation*}
        u\left(\boldsymbol{x}^{*\prime}\right) \geq u\left(\boldsymbol{x}^{\prime}\right), 
            ~ \forall \boldsymbol{x}\in \pi\left(x_1^*\right)
    \end{equation*}
    
    从而$u\left(\boldsymbol{x}^{*\prime}\right)=v\left(x_1^*\right)$，带入上式就得到了$t=0$时定理的结论。其余可依此类推即可。
\end{proof}

\cref{th-policy-sp-fe}表明，$SP$问题的最优计划可以构成$FE$的决策规则。
下面一个定理则提供了相反方向的直觉：它表明，在一定的有界条件下，符合\cref{eq-opt-plan-sp}的序列是一个最优计划。

\begin{theorem} \label{th-policy-fe-sp}
    令$X, \Gamma, F, \beta$符合\cref{as-nonempty-plan}和\cref{as-lim-exist}。
    令$\boldsymbol{x}^*\in \pi\left(x_0\right)$，且满足：
    \begin{equation*}
        v^*\left(x_t^*\right) = F\left(x_t^*,x_{t+1}^*\right) + \beta v^*\left(x_{t+1}^*\right), 
            \quad t=0,1,2,\cdots
    \end{equation*}
    并且有$\lim\limits_{t\rightarrow\infty}\sup\beta^t v^*\left(x_t^*\right)\leq 0$成立。
    那么，$\boldsymbol{x}^*$达到$SP$的上确界，即是最优计划。
\end{theorem}
\begin{proof}
    由条件：
    \begin{equation*}
        \begin{aligned}
            v^*\left(x_0\right) = u\left(\boldsymbol{x}^*\right)&= F\left(x_0^*,x_1^*\right)
                +\beta v^*\left(x_1^*\right) \\
            &= \cdots \\
            &= u_n\left(\boldsymbol{x}^*\right) + \beta^{n+1}v^*\left(x_{n+1}^*\right)
        \end{aligned}
    \end{equation*}
    又由有界条件可知：$v^*\left(x_0\right)\leq u\left(\boldsymbol{x}^*\right)$。
    由于$\boldsymbol{x}^*\in \pi\left(x_0\right)$：$v^*\left(x_0\right)\geq u\left(\boldsymbol{x}^*\right)$
    从而：$v^*\left(x_0\right) = u\left(\boldsymbol{x}^*\right)$，由定义得证。
\end{proof}

\begin{definition}
    如果对任意非空对应$G: X\rightarrow X$，成立$G\left(x\right)\subseteq \Gamma\left(X\right)$，
    则称$G$为策略对应（Policy Correspondence）。若$G$为单值的，我们称其为策略函数（Policy Function），记为$g$。
    如存在序列$\boldsymbol{x}=\left(x_0,x_1,\cdots\right)$满足
    $x_{t+1}\in G\left(x_t\right), t=0,1,2,\cdots$，则称$\boldsymbol{x}$由$G$从$x_0$生成。
    可以定义最优策略对应（Optimal Policy Correspondence）$G^*$：
    \begin{equation*}
        G^*(x)=\left\{y \in \Gamma(x): v^*(x)=F(x, y)+\beta v^*(y)\right\}
    \end{equation*}
\end{definition}

\cref{th-policy-sp-fe}表明：每个最优计划$\{x_t^*\}$都由$G^*$生成；
\cref{th-policy-fe-sp}表明：每个由$G^*$生成的计划$\{x_t^*\}$，若满足有界条件，就是一个最优计划。


% 2.4.3 FE解的存在性
\subsection{最大化定理和不动点的存在与唯一性}

我们在\ref{sec-math-pre}节最后提出，对于递归形式的最优增长问题，我们希望利用Blackwell定理证明$T$是压缩映射，
进而运用压缩映射定理证明贝尔曼方程$v=Tv$这一不动点问题解的存在和唯一性。我们已经证明了贝尔曼算子$T$是压缩映射，
但是在运用压缩映射定理时遇到了困难：我们希望值函数处于连续空间，确保空间的完备性，这样才能使用压缩映射定理。
在\ref{sec-value-func-property}节，我们将动态规划写成了更一般化的形式，只要完成一般化形式下的证明，
那么RCK模型中相关结论自然成立。我们还需要一些额外假设和定理的辅助。
方便起见，在此处写下我们所需研究的一般化形式泛函方程：
\begin{equation}    \label{eq-general-fe}
    v\left(x\right)=\max_{y\in\Gamma\left(x\right)}\left[F\left(x,y\right)+\beta v\left(y\right)\right]
\end{equation}

\begin{assumption}[凸性假设] \label{as-convex-X}
    状态变量所有可能取值集合$X$是$\mathbb{R}^l$的凸子集，可行性约束对应$\Gamma :X\rightarrow X$是非空、连续、紧集。
\end{assumption}

\begin{assumption}[有界回报]  \label{as-bounded-returns}
    函数$F: A\rightarrow \mathbb{R}$是有界连续的，且$0<\beta<1$。
\end{assumption}

显然，\cref{as-convex-X}和\cref{as-bounded-returns}蕴含了\cref{as-nonempty-plan}和\cref{as-lim-exist}。
为了应用压缩映射定理，我们不仅希望值函数有界，更希望在完备空间中研究值函数，因此
一个自然的想法就是定义了上确界范数$\|f\|=\sup _{x\in X}|f\left(x\right)|$的有界连续函数
$f: X\rightarrow \mathbb{R}$构成空间$C\left(X\right)$。定义空间$C\left(X\right)$上的算子$T$：
\begin{equation}
    Tf\left(x\right)=\max_{y\in\Gamma\left(x\right)}\left[F\left(x,y\right)+\beta f\left(y\right)\right]
\end{equation}
由此，我们将\cref{eq-general-fe}简化为不动点问题$v=Tv$。
我们将借助最大化定理（Maximum Theorem）（\cref{th-max-th}）、Blackwell充分条件（\cref{th-blackwell}）和
压缩映射定理（\cref{th-cmt}）证明：算子$T$将$C\left(X\right)$中的元素（函数）映射到$C\left(X\right)$上，
即$T: C\left(X\right)\rightarrow C\left(X\right)$，从而$T$在有界连续函数空间$C\left(X\right)$上存在唯一不动点。

\begin{theorem}[Maximum Theorem]    \label{th-max-th}
    令$X\subseteq\mathbb{R}^l$，$Y\subseteq\mathbb{R}^m$，令$f: X\times Y \rightarrow \mathbb{R}$是连续函数，
    令$\Gamma: X\rightarrow Y$是紧值连续对应（Correspondence），那么：
    \begin{itemize}
        \item 从$X\rightarrow R$的函数$h\left(x\right)\equiv\max _{y\in\Gamma\left(x\right)}f\left(x,y\right)$是连续的
        \item 由$G(x)\equiv\underset{y\in\Gamma(x)}{\operatorname{argmax}}f(x, y)=\{y\in\Gamma(x): f(x, y)=h(x)\}$
        给定的解对应$G: X\rightarrow Y$是非空、紧值和上半连续的。
    \end{itemize}
\end{theorem}

对于形如：$h\left(x\right)\equiv \max_{y\in\Gamma\left(x\right)}f\left(x,y\right)$，
最大化定理给出了$h\left(x\right)$和相应最优$y$值的集合$G\left(x\right)$随$x$连续变化的条件。
同时，若$f\left(x,y\right)$关于$y$严格凹，并且$\Gamma\left(x\right)$是凸的，那么$G$是单值的，可记为函数$g\left(x\right)$。

\begin{theorem}[$FE$问题解的存在及唯一性] \label{th-fe-solution-exist}
    令$X, \Gamma, F, \beta$满足\cref{as-convex-X}和\cref{as-bounded-returns}，
    $C\left(X\right)$为定义了上确界范数的有界连续函数空间，那么：
    \begin{itemize}
        \item 算子$T$将$C\left(X\right)$映射到自身
        \item $T$有唯一不动点$v\in C\left(X\right)$
        \item 任意给定$v_0\in C\left(X\right)$，成立：$\|T^nv_0-v\|\leq \beta^n \|v_0-v\|$
    \end{itemize}
\end{theorem}

\begin{proof}
    \textbf{Step1：证明算子$T$将有界连续函数空间映射到自身}

    注意到在\cref{as-convex-X}和\cref{as-bounded-returns}下，$F$和$f$均有界连续，集合$\Gamma$是非空、紧值和连续的。
    因此，对任意$f\in C\left(X\right), x\in X$，如下问题：
    \begin{equation*}
        Tf\left(x\right)=\max_{y\in\Gamma\left(x\right)}\left[F\left(x,y\right)+\beta f\left(y\right)\right]
    \end{equation*}
    等价于在紧集$\Gamma\left(x\right)$上最大化一个连续函数，从而最大值必能取到。同时，由于$F$和$f$均有界，从而$Tf$有界。
    最后，由最大化定理（\cref{th-max-th}），$Tf$（即定理中的函数$h$）连续。
    从而证明了：$T: C\left(X\right)\rightarrow C\left(X\right)$。

    \textbf{Step2：利用Blackwell条件，证明算子$T$是压缩映射}

    Blakwell充分条件（\cref{th-blackwell}）的前提：定义了上确界范数的有界函数空间，显然$C\left(X\right)$符合。
    接着验证两个条件:

    （1）单调性：

    令$f,h \in C\left(X\right), f\leq h$，对$f$应用算子$T$：
    $$Tf\left(x\right)=\max_{y\in \Gamma\left(x\right)}\left[F\left(x,y\right)+\beta f\left(y\right)\right]$$
    对$\forall x\in X$，令$g\left(x;f\right)$为达到最大值的$y$的集合：
    \begin{equation*}
        g\left(x;f\right) \equiv \{y\in\Gamma\left(x\right)|Tf\left(x\right)
            =F\left(x,y\right)+\beta f\left(y\right)\}
    \end{equation*}
    由于：$f\left(y\right)\leq h\left(y\right), \forall y\in \Gamma\left(x\right)$
    从而：
    \begin{equation*}
        Tf\left(x\right) = F\left(x,g\left(x;f\right)\right)+\beta f\left(g\left(x;f\right)\right)
            \leq F\left(x,g\left(x;f\right)\right)+\beta h\left(g\left(x;F\right)\right), ~\forall x\in X
    \end{equation*}
    注意到：
    \begin{equation*}
        F(x, g(x ; f))+\beta h(g(x ; f)) \leq 
            \max _{y \in \Gamma(x)}[F(x, y)+\beta h(y)]=T h(x), ~\forall x \in X
    \end{equation*}
    因此有：$Tf\left(x\right)\leq Th\left(x\right), \quad x\in X$，即$T$满足单调性。
    
    （2）折扣性：

    对任意$f\in C\left(X\right)$：
    \begin{equation*}
        \begin{aligned}
            T\left(f+a\right)\left(x\right) &= 
                \max_{y\in\Gamma\left(x\right)}\left[F\left(x,y\right) + 
                    \beta \left(f\left(y\right)+a\right)\right]  \\
            &= \max_{y\in\Gamma\left(x\right)}\left[F\left(x,y\right) + 
                \beta \left(f\left(y\right)\right)\right]+\beta a
        \end{aligned}
    \end{equation*}
    即：$T\left(f+a\right)\left(x\right)=T\left(x\right)+\beta a, \quad f\in C\left(X\right)$。

    因此，算子$T$符合Blackwell充分条件，从而是$C\left(X\right)$上的压缩映射。

    \textbf{Step3：应用压缩映射定理}

    我们已经证明，定义了上确界范数的有界连续函数空间$C(X)$是完备空间。因此由压缩映射定理（定理~\ref{th-cmt}）立得后两条结论。
\end{proof}

\cref{th-fe-solution-exist}是说明值函数存在和唯一性的关键一步，它表明：

（1）对于形如\cref{eq-general-fe}的泛函问题（$FE$）
$v\left(x\right)=\max_{y\in\Gamma\left(x\right)}\left[F\left(x,y\right)+\beta v\left(y\right)\right]$：
不动点存在且唯一；

（2）数值解法的合理性：给定一个初始猜测，通过迭代算子$T$可以收敛到不动点。

回顾\cref{th-solution-fe-sp}：满足特定有界条件的$FE$问题的解$v$就是$SP$问题的解$v^*$（上确界函数）。
根据这一最优性原理：在\cref{th-fe-solution-exist}的假设下，满足一般化$FE$问题\cref{eq-general-fe}：
$v\left(x_t\right)=F\left(x_t, x_{t+1}\right)+\beta v\left(x_{t+1}\right), ~ t=0,1,2, \ldots$
的唯一有界连续函数$v$就是对应序列问题（$SP$）的上确界函数。也就是说，在\cref{as-convex-X}和\cref{as-bounded-returns}下，
\cref{th-solution-fe-sp}和\cref{th-fe-solution-exist}共同证明：\textbf{上确界函数有界并且连续}。

回顾\cref{th-policy-fe-sp}：满足
$v^*\left(x_t^*\right) = F\left(x_t^*,x_{t+1}^*\right) + \beta v^*\left(x_{t+1}^*\right), ~ t=0,1,2,\cdots$
和特定有界性条件的任意序列是一个最优计划。因此，结合\cref{th-policy-fe-sp}和\cref{th-fe-solution-exist}可知，
至少存在一个最优计划：由非空对应$G$产生的任意计划是最优的。


% 2.4.4 ===== 值函数性质 =====
\subsection{值函数与策略函数的性质}

在相关规范化假设下（如$F$的单调性、凹性等），结合上面的论证，我们可以总结出如下性质：
\begin{property}
    值函数$v$存在且唯一；任给初始函数，通过迭代贝尔曼算子，可以收敛至值函数$v$。
\end{property}
\begin{proof}
    由\cref{th-fe-solution-exist}立得。
\end{proof}

\begin{lemma}
    令$(S,d)$为完备度量空间，令$T:S\rightarrow S$是不动点$v\in S$，那么：
    \begin{itemize}
        \item 若$S^{\prime}$是$S$的闭子集，且$T(S^{\prime})\subseteq S^{\prime}$，
            则有$v\in S^{\prime}$；
        \item 若$T(S^{\prime})\subseteq S^{\prime\prime} \subseteq S^{\prime}$，则有$v\in S^{\prime\prime}$。
    \end{itemize}
\end{lemma}
\begin{property}
    值函数$v$严格递增、严格凹并且对应的决策规则$g$是单值函数。
\end{property}
\begin{proof}
    这两条结论均可由压缩映射的性质得到。直观的证明思路是：
    假设初始函数$v_0$具有严格递增且严格凹的性质，并且在迭代过程中贝尔曼算子保留这些性质，从而极限（不动点处）也会保留这些性质。

    先证明单调性：对贝尔曼方程：
    \begin{equation*}
        v(x)=\max _{y \in \Gamma(x)}[F(x, y)+\beta v(y)], ~\forall x \in X
    \end{equation*}
    令$f\in C(X)$，令$x<x^{\prime}$，从而有：
    \begin{equation*}
        \begin{aligned}
            T f(x) & =\max _{y \in \Gamma(x)}[F(x, y)+\beta f(y)] \\
            & \leq \max _{y \in \Gamma\left(x^{\prime}\right)}[F(x, y)+\beta f(y)] \\
            & < \max _{y \in \Gamma\left(x^{\prime}\right)}\left[F\left(x^{\prime}, y\right) + 
                \beta f(y)\right]=T f\left(x^{\prime}\right)
        \end{aligned}
    \end{equation*}
    也就是说，算子$T$将有界连续函数映射到严格增的函数空间上。记$C^{\prime}(X)$为$X$上的有界连续非减函数空间，
    令$C^{\prime\prime}(X)\subset C^{\prime}(X)$是有界连续严格增函数空间。
    那么，由于$T$将$f\in C^{\prime}(X)$映射到$Tf\in C^{\prime\prime}(X)$，从而不动点$v=Tv\in C^{\prime\prime}(X)$，
    即值函数$v$严格增。

    下面证明凹性：令$x \neq x^{\prime}$，$x_\theta=\theta x+(1-\theta) x^{\prime}, \theta \in(0,1)$；
    令$y \in \Gamma(x)$使得$T f(x)=F(x, y)+\beta f(x)$；类似的，令$y^{\prime} \in \Gamma\left(x^{\prime}\right)$使得
    $T f\left(x^{\prime}\right)=F\left(x^{\prime}, y^{\prime}\right)+\beta f\left(x^{\prime}\right)$，
    从而有：
    \begin{equation*}
        \begin{aligned}
            T f\left(x_\theta\right) & \geq F\left(x_\theta, y_\theta\right)+\beta f\left(y_\theta\right) \\
            & >\theta[F(x, y)+\beta f(y)]+(1-\theta)\left[F\left(x^{\prime}, y^{\prime}\right)+
                \beta f\left(y^{\prime}\right)\right] \\
            & =\theta T f(x)+(1-\theta) T f\left(x^{\prime}\right)
        \end{aligned}
    \end{equation*}
    即$Tf$严格凹。与上面类似，记$C^{\prime}(X)$为有界连续凹函数空间，令$C^{\prime\prime}(X)\subset C^{\prime}(X)$
    是有界连续严格凹函数空间。那么，由于$T$将$f\in C^{\prime}(X)$映射到$Tf\in C^{\prime\prime}(X)$，
    从而不动点$v=Tv\in C^{\prime\prime}(X)$，即值函数$v$严格凹。

    又由于$F$关于$y$严格凹，可行集$\Gamma(x)$为凸集，根据最大化定理可知存在唯一$y=g(x)$达到最大值。
\end{proof}

\begin{property}
    策略函数$g$单调递增。
\end{property}
\begin{proof}
    由一阶条件\cref{eq-bellman-foc}：
    \begin{equation*}
        -\mathcal{F}_2(x,x^{\prime}) = \beta v^{\prime}(x^{\prime})
    \end{equation*}
    由于$\mathcal{F}(x,x^{\prime})$关于第二个变量严格凹，因此左侧关于$x^{\prime}$递增；
    由于$v$时凹的，因此右侧关于$x^{\prime}$严格递减。
    又由于右侧独立于$x$，左侧关于$x$递增，因此决策规则$x^{\prime}=g(x)$关于$x$递增。
\end{proof}